\section{第八天}
\begin{enumerate}
  \item 法治在社会经济发展当中的作用。
        \begin{enumerate}
          \item 以法治保障\zx{经济发展}。
          \item 以法治保障\zx{政治稳定}。
          \item 以法治保障\zx{文化繁荣}。
          \item 以法治保障\zx{社会和谐}。
          \item 以法治保障\zx{生态良好}。
          \item 备考提示。
                发挥法治在经济社会发展中的作用这部分内容相对比较空泛，基本的答题思路是在\bs{五位一体}的社会发展布局当中，
                在实现\zx{物质文明、政治文明、精神文明、社会文明和生态文明}的社会发展目标的过程里都\bs{离不开法治的保障作用}，
                如果在题目当中问到这个问题，从这五个方面展开即可，
                每一个点说点相关的东西，有一个 3 - 5 句话的篇幅就可以了。
                典型如：
                \begin{enumerate}
                  \item 以法治保障经济发展，\bs{运用法治方式保障}市场良性竞争，维护稳定的市场秩序，不断推动市场开放，营造良好的营商环境；
                  \item 以法治保障政治稳定，\bs{进一步}完善民主政治相关的法律制度，保障人民当家作主的地位，扩大公民民主政治权利的范围，通过法治方式巩固人民民主专政的国体和人民代表大会的政体，消除影响政治稳定的各种因素。
                \end{enumerate}
        \end{enumerate}
  \item 正确处理全面依法治国的重大关系。
        \begin{enumerate}
          \item 法治与政治。
                \begin{enumerate}
                  \item 党大还是法大是一个\zx{伪命题}，
                        依法治国、党的领导和人民当家做主是\zx{有机统一}的，
                        \zx{不存在对立关系}。
                  \item \zx{法大还是权大}才是一个真命题，
                  \item 各级领导干部必须牢记\zx{法律底线}不可碰触，\zx{破除人治}思维，
                  \item 禁止任何个人和组织\zx{在法律之上}行使权力，
                  \item 坚决避免\zx{以言代法，以权压法，徇私枉法，逐利违法}现象的发生。
                \end{enumerate}
          \item 法治与\bs{改革}。
                \begin{enumerate}
                  \item 法治\bs{和改革}的总体关系。
                        \begin{enumerate}
                          \item 法治和改革有着\zx{内在的必然联系}，
                                二者\zx{相辅相成、相伴而生，如鸟之两翼、车之两轮}。
                          \item 必须\zx{在法治下推进}改革，\zx{在改革中完善}法治。
                          \item 要把法治改革\zd{纳入}全面深化改革的\zx{总体部署}，
                                发挥法治对改革的\zx{引领和推动}作用，
                                确保重大改革\zx{于法有据}，
                          \item 做到\zx{在法治的轨道上推进改革}，
                          \item 以法治\zx{凝聚}改革\zx{共识}、
                                以法治\zx{引领}改革\zx{方向}、
                                以法治\zx{规范}改革\zx{进程}、
                                以法治\zx{化解}改革\zx{风险}、
                                以法治\zx{巩固}改革\zx{成果}。
                        \end{enumerate}
                  \item 立法决策与改革决策。
                        \begin{enumerate}
                          \item 要坚持改革决策和立法决策\zx{相统一、相衔接}，确保改革和法治实现\zx{良性互动}。
                          \item 要发挥立法对改革的\zx{引领和推动}作用，确保重大改革\zx{于法有据}。
                          \item 对\zx{实践证明}已经\zx{比较成熟}的改革经验和\zx{行之有效}的改革举措，要尽快上升为法律，\bs{通过立法扫清改革障碍}。
                          \item 立法条件\zx{不够成熟}，需要\zx{先行试点}的，要\zx{按照法定程序}作出授权，
                          \item 既不允许改革\zx{突破法律红线}，也不允许简单以现行法律没有依据\zx{迟滞改革}。
                        \end{enumerate}
                \end{enumerate}
          \item 依法治国\bs{与以德治国}。
                \begin{enumerate}
                  \item 古代社会更加\zx{强调道德}在社会治理中的\bs{核心作用}，现代社会更加\zx{强调法律}在社会治理中的\bs{重要价值}。
                  \item 依法治国属于\zx{政治文明}范畴，以德治国属于\zx{精神文明}范畴，
                  \item 依法治国与以德治国\zx{并非彼此对立}，而是\zx{相互补充、相互促进}。
                  \item \bs{法律是成文的道德，道德是内心的法律}，
                  \item 法律和道德都具有\zx{规范社会行为、维护社会秩序}的作用。
                  \item 因此，治理国家、治理社会必须\zx{一手抓法治、一手抓德治}，
                  \item 既重视发挥法律的\zx{规范作用}，以法律增加道德的\zx{客观性和理性}，
                  \item 又重视发挥道德的\zx{教化作用}，以道德增加法律的\zx{德性}，
                  \item 实现法律和道德\zx{相辅相成}、法治和德治\zx{相得益彰}。
                  \item 依法治国和以德治国相结合的理论是对我国古代\bs{治国理政经验的传承}，
                  \item 也是对国家\zx{治理现代化理论的丰富和发展}。
                \end{enumerate}
          \item 依法治国\bs{与依规治党}。
                \begin{enumerate}
                  \item 国有国法，党有党规。
                  \item 依法治国、依法执政，
                        既要求党依据宪法法律治国理政，
                        也要求党依据党内法规管党、治党。
                  \item 依规管党、治党是依法治国的重要前提和政治保障，
                        只有把党建设好，国家才能治理好。
                  \item 要坚持依法治国与制度治党、依规治党统筹推进、一体建设，
                        注重党内法规同国家法律法规的衔接和协调。
                  \item 统筹推进依规治党和依法治国，
                        促进党的制度优势与国家制度优势相互转化。
                \end{enumerate}
        \end{enumerate}
\end{enumerate}